\section{Contraction on ordered intervals}
\label{sec:ordered-contraction}

A tuple of residue classes
\(\boldsymbol s=(s_v)_{v\in V}\in(\Z/M_0\Z)^V\) is called
\emph{good} if
\[
 s_{v_{j+1}}-s_{v_j}
 \quad\text{is a strict square class modulo }M_0
 \qquad(0\le j<24),
\]
where \(v_{24}=v_0\). Thus each difference is \(1\) modulo \(8\)
and a nonzero square at every odd prime dividing \(M_0\). Let
\begin{equation}
 \rho=M_0^{-d}\bigl|\{\boldsymbol s\in(\Z/M_0\Z)^V:
                         \boldsymbol s\text{ is good}\}\bigr|.
 \label{eq:good-residue-proportion}
\end{equation}
We will bound the contribution of each good tuple when every prime uses
the nonexceptional part of its measure. First we show that these assignments
have a fixed positive proportion; the proof also explains the cycle length
$24$. We retain the parameters and lifts from the preceding section.

\begin{lemma}[Existence of good residue assignments]
\label{lem:good-residues}
The proportion $\rho$ defined in~\eqref{eq:good-residue-proportion} satisfies
\[
 \rho\ge M_0^{-23}>0.
\]
\end{lemma}

\begin{proof}
We construct a length-$24$ sequence of admissible increments with sum zero at
each factor of $M_0$.  Modulo $8$, take all the increments equal to $1$; their
sum is $24\equiv0\pmod8$.

Let $p<P$ be an odd prime.  If $-1$ is a square modulo $p$, repeat the pair
$(1,-1)$ twelve times.  Both entries are nonzero squares and each pair has sum
zero.  If $-1$ is not a square, put
\[
 S=\{x^2:x\in\F_p\}.
\]
The sets $S$ and $-1-S$ each have $(p+1)/2$ elements, so they intersect.
Thus $-1=a+b$ for some $a,b\in S$.  Neither $a$ nor $b$ is zero, since otherwise
$-1$ would be a square.  The triple $(a,b,1)$ consists of nonzero squares and
has sum zero.  Repeat this triple eight times.  This also covers $p=3$, where
the triple is $(1,1,1)$.

The Chinese remainder theorem now gives increments
$\Delta_0,\ldots,\Delta_{23}\pmod{M_0}$, each a strict square class, with
$\sum_{j=0}^{23}\Delta_j=0$.  For each initial residue $a\pmod{M_0}$, prescribe
\[
 s_{v_0}=a,\qquad
 s_{v_j}=a+\sum_{i=0}^{j-1}\Delta_i\quad(1\le j<24).
\]
The last edge has increment $\Delta_{23}$ because the sum of all the increments is
zero.  The slots $v_0,\ldots,v_{23}$ are distinct, since their block words are
the distinct cyclic shifts of the identity word.  Repetition of residue
values at nonadjacent slots causes no difficulty.  The other $d-24$ residues
are unrestricted.  We have therefore exhibited at least
$M_0\cdot M_0^{d-24}$ good assignments among the $M_0^d$ possible assignments,
which proves the bound.
\end{proof}

The length $24$ is divisible by $8$, $2$, and $3$, exactly the three
local cycle lengths used in this construction. Thus this choice works
uniformly for every prime included in the fixed modulus.

We now prove the analytic estimate for each good assignment.

\begin{proposition}[Contraction for good residue assignments]
\label{prop:nonexceptional-contraction}
There are absolute constants \(C\) and \(N_0\) such that, whenever
\(N\ge N_0\), \(A\subseteq[N]\) has no positive square difference, and
\(\boldsymbol s\) is good, one has
\begin{equation}
 \left|\int\prod_{v\in V}
  \mathcal L_{[N],\mathcal P_N,Q}f_{s_v}(z_v)
  \,d\bigotimes_{p\in\mathcal P_N}(\mu_p-\varepsilon_p)\right|
 \le C N^{-\sigma_0}.
 \label{eq:nonexceptional-contraction}
\end{equation}
The constants are uniform in \(A\) and \(\boldsymbol s\).
\end{proposition}

\begin{proof}
Write \(\eta_p\) for the total mass of \(\varepsilon_p\), and set
\[
 \mu_p^\circ=\frac{\mu_p-\varepsilon_p}{1-\eta_p},
 \qquad \Lambda^\circ=\bigotimes_{p\in\mathcal P_N}\mu_p^\circ.
\]
It suffices to prove the assertion with \(\Lambda^\circ\) in place of the
measure in~\eqref{eq:nonexceptional-contraction}, since the latter is
\(\prod_p(1-\eta_p)\Lambda^\circ\) and \(0\le\prod_p(1-\eta_p)\le1\).
The single-slot marginals of \(\Lambda^\circ\) are uniform. Moreover, at
each prime its marginal on any directed cycle edge is uniform on strict
square pairs. We use conductor decay for \(\Lambda^\circ\), without any
claim that this normalized measure is reflection positive.

\textit{Selecting an ordered pair.}\quad
Put \(K=\lfloor N^\tau\rfloor\), and partition \([N]\) into ordered
consecutive intervals \(I_1,\ldots,I_K\) of sizes differing by at most one.
Writing \(w_i=|I_i|/N\), we have the exact identity
\[
 \mathcal L_{[N],\mathcal P_N,Q}f_s
   =\sum_{i=1}^K w_i\mathcal L_{I_i,\mathcal P_N,Q}f_s.
\]
Thus the integral is an average over independent interval choices in all
slots, with probabilities \(w_i\). For large \(N\), every interval has
length at least \(\tfrac12N^{1-\tau}\). Since \(1-\tau>6\beta\), this
length is at least \(10Q^6\). Lemma~\ref{lem:lift-moments} therefore
applies uniformly to every interval lift and every exact-denominator
truncation used below. Its application to \(f_s\), which is bounded by
the fixed constant \(M_0\), changes only the constants. All occurrences
of \(N^{o(1)}\) in this proof are uniform over these interval choices,
residues and sets.

First consider assignments for which all 24 cycle slots choose the same
interval. Their total probability is
\[
 \sum_{i=1}^K w_i^{24}
 \le(\max_i w_i)^{23}\ll N^{-23\tau}.
\]
For every assignment, ordinary H\"older and the uniform single-slot
marginals give
\[
 \left|\int\prod_v G_v(z_v)\,d\Lambda^\circ\right|
 \le\prod_v\|G_v\|_d=N^{o(1)},
 \qquad
 G_v=\mathcal L_{I(v),\mathcal P_N,Q}f_{s_v}.
\]
Here \(d\) is even, so the required moment is supplied directly by
Lemma~\ref{lem:lift-moments}. The contribution of these assignments is
therefore \(O(N^{-23\tau+o(1)})\); there is no restriction on the interval
choices in the other slots.

In every remaining assignment there is a directed cycle edge \(u\to v\)
whose first interval precedes its second. Indeed, if the cyclic sequence
of interval indices had no ascent, it would be nonincreasing all the way
around the cycle, and hence constant. Fix one ascending edge, and denote
its two intervals by \(I_1,I_2\) for the rest of this argument. All
integers in \(I_1\) precede all integers in \(I_2\).

\textit{Eliminating the other slots.}\quad
Retain these two lifts in full. In each of the other \(d-2\) slots,
truncate the lift to exact conductors at most \(T=N^\tau\).
Lemma~\ref{lem:conductor-decay}, applied during a telescoping replacement
of these slots, bounds the total error by
\begin{equation}
 O(T^{-\gamma}N^{o(1)})=O(N^{-\gamma\tau+o(1)}).
 \label{eq:interval-truncation-error}
\end{equation}
Every norm used here is controlled by Lemma~\ref{lem:lift-moments}:
the retained and discarded pieces are subfamilies of complete exact
denominator pieces. This error is estimated before the Fourier expansion
that follows.

Each Fourier coefficient in a truncated slot has modulus at most \(M_0\),
and there are at most \(\sum_{q\le T}\varphi(q)\ll T^2\) frequencies.
The product expansion of all other slots consequently has total
coefficient absolute sum at most
\begin{equation}
 C_{d,M_0}T^{2(d-2)}\ll N^{2d\tau}.
 \label{eq:interval-fourier-cost}
\end{equation}
For one term in this expansion, let \(B\subseteq\mathcal P_N\) be the
union of the prime supports of its characters, and put
\[
 R=\prod_{p\in B}p.
\]
The conductors are squarefree, so their least common multiple satisfies
\begin{equation}
 R\le T^{d-2}\le N^{d\tau}.
 \label{eq:interval-mode-modulus}
\end{equation}

Condition on all slot labels at the primes in \(B\). The expanded
characters become constants of modulus one. Since \(\Lambda^\circ\) is
a product over primes, the tuple laws at the remaining primes are
unchanged and independent. The two selected labels now have fixed
coordinates \(a_{1,B},a_{2,B}\) at \(B\), whose difference is a nonzero
square at each prime in \(B\). At every remaining prime their pair
marginal is still uniform on strict square pairs.

Let \(\mathcal S=\mathcal P_N\setminus B\), let \(s_1=s_u\) and
\(s_2=s_v\), and define functions on \(X_{\mathcal S}\) by
\[
 g_i(z)=\bigl(\mathcal L_{I_i,\mathcal P_N,Q}f_{s_i}\bigr)(a_{i,B},z),
 \qquad i=1,2.
\]
With Fourier coefficients normalized for uniform measure on
\(X_{\mathcal S}\), their conditional pair integral is exactly
\begin{equation}
 \sum_{\xi\text{ on }X_{\mathcal S}}
   \widehat g_1(-\xi)\widehat g_2(\xi)\mathfrak g(\xi).
 \label{eq:interval-pair-fourier}
\end{equation}
To see this, simultaneous translation kills character pairs whose
frequencies do not sum to zero. For a matching pair, the multiplier is
the character average of the difference of the labels. At every odd
prime a uniform nonzero square is the square of a uniform nonzero
unit. The Chinese remainder theorem therefore identifies this average
with the unit-square multiplier \(\mathfrak g(\xi)\) used in
Proposition~\ref{prop:square-kernel}.

\textit{Estimating the selected pair.}\quad
Specialization on a uniform coordinate block of size \(R\) costs at
most \(R^{1/2}\) in \(L^2\): for every function \(G\),
\[
 \|G(a_B,\cdot)\|_2^2\le R\|G\|_2^2.
\]
This inequality holds at every specified fiber, independently of the
law under which that fiber was selected. Hence
\(\|g_i\|_2\le R^{1/2}N^{o(1)}\). By the bound
\(|\mathfrak g(\xi)|\le q(\xi)^{-1/3}\), Parseval and
Cauchy--Schwarz, the part of~\eqref{eq:interval-pair-fourier} with
conductor exceeding \(H\) has modulus at most
\begin{equation}
 H^{-1/3}\|g_1\|_2\|g_2\|_2
 \le H^{-1/3}R N^{o(1)}.
 \label{eq:interval-pair-tail}
\end{equation}

We next identify the remaining coefficients with actual integer fibers.
Put \(D=M_0R\). The parameter inequalities
\begin{equation}
 d\tau<\frac{\beta\gamma}{1000}<\frac\beta2
 \label{eq:interval-parameter-room}
\end{equation}
imply \(RH\le Q\) and, for sufficiently large \(N\), \(D\le N^\beta\).
For a remaining frequency \(\xi\) with \(q(\xi)\le H\), all frequencies
\(\xi+\eta\), where \(\eta\) ranges over the character group on \(B\),
therefore occur in the original lift. Summing this entire character group
gives
\begin{align}
 \widehat g_i(\xi)
 &=\sum_{\eta\text{ on }B}\frac1{|I_i|}
     \sum_{n\in I_i}f_{s_i}(n)e(-n(\xi+\eta))e(a_{i,B}\eta)\notag\\
 &=\frac{D}{|I_i|}
     \sum_{\substack{n\in I_i\\n\equiv a_i\pmod D}}
           1_A(n)e(-n\xi),
 \label{eq:interval-selected-fiber}
\end{align}
where \(a_i\pmod D\) combines \(s_i\pmod {M_0}\) and the coordinates
\(a_{i,B}\) by the Chinese remainder theorem. Character orthogonality
supplies the factor \(R\), while the definition of \(f_{s_i}\) supplies
\(M_0\). Goodness of \(\boldsymbol s\) and the constraints at the
conditioned primes show that \(a_2-a_1\) is a strict square class
modulo \(D\).

The remaining frequencies of conductor at most \(H\) are precisely the
frequencies with squarefree conductor at most \(H\) coprime to \(D\).
Indeed, \(D\) contains every prime less than \(P\); every prime divisor
of such a conductor is consequently an available prime at least \(P\),
at most \(H\le Q\), and outside \(B\). This includes the zero
frequency under the convention \(q(0)=1\).

Define functions on \([N]\) by
\[
 F_i(n)=1_A(n)1_{I_i}(n)1_{n\equiv a_i\pmod D}.
\]
They are bounded by one, their supports are ordered, and there is no
positive square difference from the first support to the second.
Moreover, \(D\) is eight times an odd squarefree number and
\(D\le N^\beta\). All hypotheses of
Proposition~\ref{prop:square-kernel} hold. If the Fourier coefficients
of \(F_i\) use normalization \(N^{-1}\), then
\eqref{eq:interval-selected-fiber} gives
\(\widehat g_i(\xi)=DN\widehat F_i(\xi)/|I_i|\). Thus the part of
\eqref{eq:interval-pair-fourier} of conductor at most \(H\) is bounded
in modulus by
\begin{align}
 &\frac{D^2N^2}{|I_1||I_2|}
 \left|\sum_{\substack{q(\xi)\le H,\ q(\xi)\text{ squarefree}\\
                       (q(\xi),D)=1}}
      \mathfrak g(\xi)\widehat F_1(-\xi)\widehat F_2(\xi)\right|
 \notag\\
 &\hspace{2cm}\ll H^{-1/3}D\frac{N^2}{|I_1||I_2|}
 \ll H^{-1/3}D N^{2\tau}.
 \label{eq:interval-kernel-cost}
\end{align}
The factor \(D\) in Proposition~\ref{prop:square-kernel} has canceled
one of the two fiber-normalization factors.

Combining~\eqref{eq:interval-pair-tail}
and~\eqref{eq:interval-kernel-cost} bounds each conditional pair integral
by \(O(N^{-\beta/6+(d+2)\tau+o(1)})\). This is uniform over all
conditionings of positive probability. Averaging them, and then summing
the Fourier expansion using~\eqref{eq:interval-fourier-cost}, bounds the
truncated integral for each unequal-cycle interval assignment by
\[
 O\bigl(N^{-\beta/6+(3d+2)\tau+o(1)}\bigr).
\]
Averaging interval assignments incurs no further loss. Including the
equal-cycle contribution and~\eqref{eq:interval-truncation-error}, the
original integral has modulus at most
\begin{equation}
 C\left(N^{-23\tau+o(1)}+N^{-\gamma\tau+o(1)}
       +N^{-\beta/6+(3d+2)\tau+o(1)}\right).
 \label{eq:interval-three-errors}
\end{equation}
Finally, \(23\tau=5888\sigma_0\), \(\gamma\tau=4\sigma_0\), and
\[
 \frac\beta6-(3d+2)\tau
 >\beta\left(\frac16-\frac3{64000}\right)
 >\frac\beta7>2\sigma_0.
\]
The finitely many uniform moment losses in each term can be bounded by
\(N^{\sigma_0}\), after increasing the absolute threshold in \(N\).
Equation~\eqref{eq:interval-three-errors} is consequently
\(O(N^{-\sigma_0})\), as required. All estimates used absolute values
of signed Fourier lifts; no pointwise positivity of a lift was needed.
\end{proof}
